\chapter{总结与展望}

\section{工作总结}

本文以AutoPT自动化渗透测试框架为基础，围绕其在流程控制、工具能力和交互方式三方面的不足做了改进。

改动最大的是PSM流程控制部分。本文设计了统一的ReAct结构化提示词模板，为Scan、Inquire、Exploit三个Agent分别定义了互不重叠的角色边界和行为约束，解决了原版各Agent输出格式混乱、越权调用工具的问题。Scan阶段被细化为端口发现、服务知识匹配、漏洞扫描和结果评估四个子阶段，配合nmap全端口扫描和services.yml服务知识库，将攻击面覆盖从单一默认端口扩展到了全部端口空间。上下文压缩机制通过工具输出即时摘要和三层结构化组装，控制了多轮ReAct循环中的Token增长。容错提取逻辑减少了Inquire Agent因格式错误而丢失PoC的情况。Check阶段将匹配域缩小为工具输出，并引入LLM判断与正则兜底校验，避免了模型推理文本被误判为利用证据的问题。

工具层面的改进集中在三处。新增的fill\_element、select\_option、wait\_for\_selector三个浏览器扩展工具补齐了Playwright原生工具包在表单交互上的缺失。RobustReActParser通过六步降级策略适配了思维链模型的非标准输出，将原版解析失败即终止的行为改造为可恢复的容错机制。ReadHTML工具的PoC提取增强使Inquire Agent能直接拿到精简后的利用代码，减少了无效的推理轮次。

用户界面方面，本文基于Flask和原生JavaScript实现了Web Console，提供控制台、漏洞库、渗透测试面板、靶机管理、测试报告和系统设置六个模块，使常见操作不再依赖命令行。

实验数据方面，基于20个真实CVE漏洞和3种大语言模型的300次独立实验中，改进版总体成功率达到63.67\%，较原版（约25\%）提升近39个百分点。简单漏洞成功率从约32\%升至76.67\%，复杂漏洞从约17\%升至50.67\%。原版60组测试中有38组通过率为0，改进版压缩到了3组。在成本方面，以GPT-4o-mini为例，单次测试成本从原版的\$0.0099降至\$0.003109，平均耗时从161.3秒降至94.2秒。失败原因分析显示，工具执行失败仅占10.1\%，是五类失败原因中最低的一项。

\section{不足与展望}

改进版在成功率和成本上都有明显进步，但实验中也暴露了几个尚未解决的问题。

首先是Agent太容易放弃。失败案例中占比最高的就是“提前放弃”（33\%），Agent碰壁两三次后就认定当前漏洞不通，主动终止了流程。Check阶段的\_build\_failure\_guidance机制虽然会根据上次失败给出重试建议，但它给的建议不区分失败类型，不管是网络超时还是命令写错还是权限不够，都是笼统地“换个方式再试”。如果能根据失败原因的具体类型分别给出针对性的重试策略，Agent坚持下去的概率应该会更高。

其次是模型的安全策略会拒绝生成攻击性命令，这在渗透测试场景下很难绕开。安全审查拦截占了失败原因的22\%，GPT-4o-mini尤为明显，达到30.77\%。要缓解这个问题，一种做法是把敏感操作拆得更碎，让每一步单独看起来不那么危险；另一种做法是换用安全限制更宽松的本地开源模型（如DeepSeek），但这需要额外评估其在、是否遵循ReAct格式和推理质量上是否够用。

最后，当前实验只覆盖了20个Web应用层面的CVE漏洞（命令注入、文件读取、反序列化等），没有涉及内网横向移动、权限提升、社会工程这些渗透测试的后续阶段。PSM的状态机结构本身支持扩展，在图中增加新的状态节点和转移规则即可，但每个新阶段都需要配套的工具、提示词和验证逻辑，工作量较大。
